\scanpage{117}
la structure envisagée peut se décrire par

a) famille d'espaces analytiques $(D_d)_{d\in\mathbb N^{\deleted{\unsure{*}}}}$

b) $\forall\,d\in\mathbb N$, un revêtement principal \inserted{$D^!_d$} de $D_d$, de groupe $\mathfrak S_d$ --- d'où des $D_{d,d'}$ $(0\leq d'\leq d)$

\margin{on aura $D_{d,d}\simeq D_{d,0}\simeq D_d$\par $D_{d,1}$ est le rev. à $d$ feuillets de $D_d$ qui correspond à $D^!_d$.}

c) des morphismes
\diagram{\begin{tikzcd}[ampersand replacement=\&,column sep=large]
D_{d,d'}\arrow[r,dash,"b",to path={-- ([xshift=-1ex]\tikztotarget.west) \tikztonodes -- ++(0,-0.07)}]
\& D_{d'} \& (\text{pour }0\leq d'<\unsure{d^\circ|d''}),
\end{tikzcd}}
et en plus des diagrammes cartésiens

\deleted{avec la condition que les morphismes can. (pour $d\unsure{\geq|>}d'\unsure{\geq|>}d''$)}
\[\deleted{D_{d,d',d''}\overset{\mathrm{def}}{=}D_{d,d'}\mathbin\times_{D_{d'}}D_{d',d''}\longrightarrow D_d\mathbin{\underset{\unsure{?}}{\times}}D_{\unsure{d''}}}\]
\deleted{d) $D_{d,d'}\mathbin\times_{D_{d'}}D^!_{d'}\unsure{\sim|\simeq}$}
\deleted{d) $\unsure{D^!_{d'}|D^!_{d,d'}}$}
\deleted{$\unsure{B_{d,d'}|D_{d,d'}}$}

d)
\diagram{\begin{tikzcd}[ampersand replacement=\&,column sep=large,row sep=large]
\&D^!_{d,d'}\arrow[r,very thick]\arrow[d,"\mathrm{can}"']\arrow[dr,phantom,"\uline{\mathrm{cart}}"]
\&D^!_{d'}\arrow[d,"\mathrm{can}"]\\
(175)\&D_{d,d'}\arrow[r,dash,"b"',to path={-- ([xshift=-1ex]\tikztotarget.west) \tikztonodes -- ++(0,0.07)}]
\&D_{d'}
\end{tikzcd}}

où $D^!_{d,d'}$ est le revêtement $\mathfrak S_{d'}$-principal de
\[D_{d,d'}\simeq D^!_d\mathbin\times_{(\mathfrak S_d)_{D_d}}\mathcal P_{d'}(I_d),\]
associé au $d'$-rev. ``tautologique'' de $D_{d,d'}$; --- ce qui donne, pour $D_{d,d',d''}$ déduit de ce rev. de $D_{d,d'}$ par
\[D_{d,d',d''}\overset{\text{déf}}{=}D^!_{d,d'}\mathbin\wedge_{(\mathfrak S_{d'})_{D_{d,\unsure{d'|d''}}}}\mathcal P_{d''}(I_{d'}),\]

\margin{\uline{NB} $D_{d,d'}$ et $D_{d,d',d''}$ sont déduits directement de $D^!_d$.}
un \deleted{\unsure{rev.}} diagramme cartésien
\diagram{\begin{tikzcd}[ampersand replacement=\&,column sep=large,row sep=large]
\&D_{\unsure{d^\circ|d''},d',d''}\arrow[r]\arrow[d]\arrow[dr,phantom,"\mathrm{cart}"]
\&D_{d',d''}\arrow[d]\\
(176)\&D_{d,d'}\arrow[r]\&D_{d'}
\end{tikzcd}}

\margin{116}
